\documentclass[12pt]{article}
\usepackage[T1]{fontenc}
\usepackage[utf8]{inputenc}
\usepackage{lmodern}
\usepackage{textcomp}
\usepackage{enumitem}
\usepackage{geometry}
\geometry{margin=1in}
\begin{document}
\begin{center}
Answer Key - Physics - Undergraduate Level Assessment 1 \
\small (Answers are ordered exactly as in the question paper.)
\end{center}

\section*{Multiple Choice}
\begin{enumerate}[label=\arabic*.]
\item \textbf{Answer:} D
\\ \textit{Options:} A) (1/2) k T; B) kT; C) (3/2) k T; D) 3 kT
\\ \textit{Note:} The average total energy of a three-dimensional harmonic oscillator in thermal equilibrium at temperature T is given by the equipartition theorem, which states that each degree of freedom contributes an average energy of (1/2)kT, resulting in a total of 3 kT for the three spatial dimensions.
\item \textbf{Answer:} A
\\ \textit{Options:} A) real; B) imaginary; C) degenerate; D) linear
\\ \textit{Note:} The eigenvalues of a Hermitian operator are always real because Hermitian operators, which represent observable physical quantities in quantum mechanics, have the property that their eigenvalues must correspond to measurable values, and the mathematical formulation ensures that these values are real numbers.
\item \textbf{Answer:} B
\\ \textit{Options:} A) an average of 10 times, with an rms deviation of about 4; B) an average of 10 times, with an rms deviation of about 3; C) an average of 10 times, with an rms deviation of about 1; D) an average of 10 times, with an rms deviation of about 0.1
\\ \textit{Note:} The quantum efficiency of 0.1 means that each photon has a 10\% chance of being detected, so with 100 photons, the expected average detection is 10, while the rms deviation, calculated using the binomial distribution, is approximately 3, reflecting the statistical fluctuations in the detection process.
\item \textbf{Answer:} A
\\ \textit{Options:} A) decreased by a factor of 9; B) decreased by a factor of 49; C) decreased by a factor of 81; D) increased by a factor of 9
\\ \textit{Note:} The emission spectrum of Li++ is identical to that of hydrogen because both have one electron, and since the energy levels are scaled by the square of the atomic number ($Z^2$), for Li++ (Z=3), the wavelengths are decreased by a factor of 9 ($3^2$) compared to hydrogen.
\item \textbf{Answer:} D
\\ \textit{Options:} A) 0.1 GeV/$c^2$; B) 0.2 GeV/$c^2$; C) 0.5 GeV/$c^2$; D) 1.0 GeV/$c^2$
\\ \textit{Note:} The rest mass can be determined using the energy-momentum relation, E² = (pc)² + (m₀c²)², where the total energy is 5.0 GeV and momentum is 4.9 GeV/c, leading to a calculated rest mass of approximately 1.0 GeV/c².
\end{enumerate}

\section*{True / False}
\begin{enumerate}[label=\arabic*.]
\setcounter{enumi}{5}
\item \textbf{Answer:} False
\\ \textit{Options:} A) True; B) False
\\ \textit{Note:} The statement is false because the time required for one revolution of a satellite, known as its orbital period, is influenced by the mass of the planet it orbits, as described by Kepler's Third Law of planetary motion, which shows that a more massive planet exerts a stronger gravitational pull, affecting the satellite's orbital dynamics.
\item \textbf{Answer:} True
\\ \textit{Options:} A) True; B) False
\\ \textit{Note:} The electric displacement current is defined by Maxwell's equations as being directly proportional to the time rate of change of the electric flux through a surface, which accounts for the displacement current in regions where electric fields vary with time.
\item \textbf{Answer:} True
\\ \textit{Options:} A) True; B) False
\\ \textit{Note:} The statement is true because the ionization energy of helium from its ground state is indeed 24.6 eV, reflecting the energy needed to remove one electron from its tightly bound 1 s orbital.
\item \textbf{Answer:} False
\\ \textit{Options:} A) True; B) False
\\ \textit{Note:} The statement is false because the relativistic momentum of a particle is given by the formula p = γmv, where γ (gamma) is the Lorentz factor and v is the particle's velocity, and simply knowing that the total energy is twice the rest energy does not directly lead to the conclusion that momentum equals mc/2.
\item \textbf{Answer:} True
\\ \textit{Options:} A) True; B) False
\\ \textit{Note:} Electromagnetic radiation emitted from a nucleus is most commonly in the form of gamma rays because gamma rays are the highest energy form of electromagnetic radiation released during nuclear decay processes, which often occur to stabilize the nucleus after a radioactive decay event.
\end{enumerate}

\section*{Long Form Response}
\begin{enumerate}[label=\arabic*.]
\setcounter{enumi}{10}
\item \textbf{Answer:} The Hall coefficient is a crucial parameter in determining the sign of charge carriers in a doped semiconductor. When a magnetic field is applied perpendicular to the current flow in the semiconductor, a voltage, known as the Hall voltage, develops across the material. The Hall coefficient (R\_H) is defined as the ratio of the induced Hall voltage to the product of the current and the magnetic field strength. A positive Hall coefficient indicates that the primary charge carriers are holes (positive charge), whereas a negative Hall coefficient signifies that electrons (negative charge) are the dominant carriers. This relationship allows for the identification of the type of doping in the semiconductor, which has significant implications for its electronic properties and applications in devices.
\\ \textit{Note:} correct because it accurately describes how the Hall coefficient, determined by measuring the Hall voltage in the presence of a magnetic field, indicates the type of charge carriers in a doped semiconductor, with a positive value signifying holes as the dominant carriers.
\item \textbf{Answer:} In a constant volume process involving an ideal gas, the relationship between internal energy and heat transfer is governed by the first law of thermodynamics, which states that the change in internal energy (ΔU) of the system is equal to the heat added to the system (Q) minus the work done by the system (W). Since the volume remains constant, no work is done (W = 0). Therefore, any heat transfer to the gas directly results in an increase in its internal energy, as expressed by the equation ΔU = Q. This means that when heat is added to the gas at constant volume, its internal energy increases, leading to a rise in temperature. Hence, in this context, the internal energy of an ideal gas is solely dependent on its temperature, which is influenced by the heat transfer occurring during the process.
\\ \textit{Note:} correct because, in a constant volume process for an ideal gas, the first law of thermodynamics simplifies to ΔU = Q, indicating that all heat added to the system directly increases its internal energy since no work is performed.
\end{enumerate}

\vfill
\noindent\textit{End of Answer Key}
\end{document}